\scanpage{353}
% top centre
\mspage{269}
% lines 1–8
pour \unsure{tte} catégorie $\deleted{\unsure{I}}\unsure{J|I}$ de card. réduit $\leq\pi$ (en particulier pour $\deleted{I}J\in\unsure{\Phi|\underline{\Phi}}$), tout foncteur $\varphi:\deleted{I}J\longrightarrow M=\operatorname{Ind}_\pi(C)$ admet une représentation \unsure{indicielle|indicible} $\Phi:J\times I\longrightarrow C$, avec $I$ \unsure{grd|gd} devant $\pi$. Ceci, \deleted{suffit} plus la stabilité de $C$ par $\varinjlim$ de type $J$, suffit à établir l'existence de la $\varinjlim$ de $\varphi$.

% lines 9–16
La réciproque provient du fait que si $C$ est karoubienne, alors son image essentielle dans $M$ n'est autre que $M_\pi$, laquelle est stable dans $M$ par \inserted{\unsure{ttes|les}} $\varinjlim$ de cardinal $\leq\pi$ représentables dans $M$. \deleted{\unsure{Elle}} Si $M$ est stable par $\varinjlim$ de type $\Phi$, il en est donc de $\unsure{\overline{\mathrm m}|\overline{\mathrm w}}$ de $M_\pi$, donc de $C$, \unsure{qed.|q.e.d.}

% lines 17–21
\uline{Corollaire} 4.11.2. Pour que $M$ soit stable par $\varinjlim$ filtrantes \inserted{\deleted{\unsure{avec|ou} \unsure{petites} \unsure{pour}} (resp. par petites $\varinjlim$)}, il suffit que $C$ le soit par $\varinjlim$ filtrantes \inserted{(resp. par $\varinjlim$)} \uline{de cardinal} $\leq\pi$, et c'est aussi nécessaire si $C$ est karoubienne.

% line 22
Supposons $C$ stable par $\varinjlim$ filtr. de
% archive sheet number: 352
